\documentclass[11pt,a4paper]{article}
\usepackage{harmonikabau_preamble}

\newcommand{\hbdoknr}{0015}
\newcommand{\hbtitel}{Acoustic Coupling and Impedance Matching}
\newcommand{\hbuntertitel}{Channel Geometry, Slot Shape and Reed Block Design}
\newcommand{\hbheadtext}{Doc.\,0015 --- Acoustic Coupling and Impedance Matching}
\newcommand{\hbfoottext}{Cross-references: Doc.\,0002 $\cdot$ Doc.\,0007 $\cdot$ Doc.\,0011}

\AtBeginDocument{%
  \fancyhead[L]{\small\hbheadtext}%
  \fancyhead[R]{\small\thepage}%
  \fancyfoot[C]{\small\color{hb@darkgray}\hbfoottext}%
}

\begin{document}
\hbmaketitle
\begin{tcolorbox}[colframe=red!60!black, colback=red!5!white]
\textbf{Note:} This document presents an explanatory model.
The numerical values given are order-of-magnitude estimates ---
they provide no constructive specifications.
The optimum is heard, not calculated.
\end{tcolorbox}

\noindent\textit{Analysis of acoustic coupling between reed and chamber as a two-port system. Transition geometry, constriction as acoustic mass, coupling centroid and impedance mismatch.
Reference: Doc.~0005--0013.}

\section{The Chamber as a Two-Port System}

In the previous documents (Doc.~0004--0008), the chamber was treated as a Helmholtz resonator with a single opening (valve). Doc.~0013 showed that the reed never completely closes the channel --- the side gap always remains open. The chamber therefore has \textbf{two ports}:

\textbf{Port~1 (reed end):} The reed pulsates and creates pressure impulses. The effective area fluctuates between the side gap ($A_\mathrm{min} \approx 10~\mathrm{mm}^2$) and the slot area ($A_\mathrm{max} \approx 26~\mathrm{mm}^2$).

\textbf{Port~2 (valve):} Stationary opening with $A_2 \approx 300~\mathrm{mm}^2$. Sound radiates from here.

The conductance $G = A/l$ determines how easily air flows through each port. Port~1 (reed) carries time-averaged $\approx 70\,\%$ of total conductance --- the reed is the primary port, not the valve.

\begin{equation}
f_H = \frac{c}{2\pi}\sqrt{\frac{S_1/l_1 + S_2/l_2}{V}}
\end{equation}

\begin{tcolorbox}[colframe=keygreen, colback=green!5!white]
\textbf{Consequence:} All calculations in Doc.~0004--0008 that derive $f_H$ only from the valve opening must be understood as lower bounds. The actual chamber resonance lies higher.
\end{tcolorbox}

\section{Is the Fold Acoustically Visible?}

For a bass at 50\,Hz, the wavelength is $\lambda = c/f = 6860$\,mm. The largest chamber dimension (chamber length 120\,mm) is only 1.7\,\% of $\lambda$. A partition wall (30\,mm) is 0.44\,\% of $\lambda$.

When $d/\lambda < 0{.}1$, the wave diffracts completely around the obstacle --- it does not exist acoustically. When $d/\lambda > 0{.}5$ it acts as a reflector.

At 50\,Hz, \textbf{all} chamber geometries lie deep in the regime $d \ll \lambda$. Only above $\approx 1400$\,Hz (28th overtone) does the partition wall begin to approach the limit $d/\lambda = 0{.}1$.

The analogy with an optical waveguide confirms this: a smooth, round tube with constant cross-section has no reflection coefficient --- the wave follows the curve as long as the bend radius is large compared to the wavelength.

\textbf{Coltman correction:} Coltman compensates not the bend itself, but the cross-section change and path-length difference at the fold. $Z = \rho c / A$ remains constant at constant cross-section. The path-length correction scales with $d^2/r$ --- at $d/\lambda < 0{.}01$ it is negligible.

\begin{tcolorbox}[colframe=keygreen, colback=green!5!white]
\textbf{Conclusion:} In the bass chamber, the fold plays no role for the fundamental and the first overtones. Matching occurs at the openings (Port~1 and Port~2), not at the fold.
\end{tcolorbox}

\section{Constriction = Acoustic Mass}

Every cross-section reduction in the channel acts as a concentrated acoustic mass: $m = \rho \times l / A$. This holds regardless of whether the constriction is in a straight tube, at a fold, or at the slot$\to$chamber transition.

The Coltman compensation in folded organ pipes is fundamentally a constriction at the deflection point. The mechanism changes (waveguide vs.\ lumped element), but the physical element (constriction = mass) remains the same.

\section{Pressure Build-up Before the Constriction}

The reed's pressure impulse backs up before the constriction at the slot$\to$chamber transition. Three scenarios:

\textbf{Too wide:} Pressure escapes immediately into the volume --- weak back-reaction, weak coupling.

\textbf{Optimal:} Pressure briefly backs up --- maximum back-reaction, best coupling.

\textbf{Too narrow:} The impulse is predominantly reflected --- energy does not enter the volume.

\section{Impedance Matching: Transition Geometry as Transformer}

The observation that different transition geometries change \textbf{response} but not \textbf{filter action} can be modelled cleanly as a matching problem:

\textbf{Source:} Bellows pressure $p_\mathrm{bellows}$ with internal resistance $R_i$ depending on the system's flow resistance. $R_i \propto 1/A^2$ (Bernoulli).

\textbf{Load:} The reed with mechanical impedance $R_L$ (fixed, determined by stiffness and mass, Doc.~0012).

\textbf{Transformer:} The transition geometry at the reed end determines how efficiently flow energy is converted to reed energy.

\begin{equation}
P / P_\mathrm{max} = \frac{4 \, R_L \, R_i}{(R_i + R_L)^2}
\end{equation}

Maximum at $R_i = R_L$ --- classical matching condition. Too narrow: $R_i \gg R_L \to$ blockage. Too wide: $R_i \ll R_L \to$ short circuit. Optimal: $R_i = R_L \to$ maximum power transfer bellows$\to$reed.

\begin{tcolorbox}[colframe=keygreen, colback=green!5!white]
\textbf{Crucial:} The filter action ($f_H$, notch filter) depends on the resonator --- chamber volume and opening area. The transformer (transition geometry) changes matching efficiency, not resonance. This is exactly what is observed: different response, same timbre.
\end{tcolorbox}

\section{Equivalent Circuit: Two Coupled Resonators}

The separation of response and filter action can be represented exactly as an electroacoustic analogy. Reed and chamber are each a series resonator (RLC), connected by a coupling element:

\textbf{Circuit 1 (reed):} Mass $m \to$ inductance $L_1$. Stiffness $k \to$ capacitance $C_1 = 1/k$. Damping $\to$ resistance $R_1$. $f_1 = 1/(2\pi\sqrt{L_1 C_1}) \approx 50$~Hz.

\textbf{Circuit 2 (chamber):} Acoustic mass $\rho l/A \to L_2$. Volume compliance $V/(\rho c^2) \to C_2$. $f_H = 1/(2\pi\sqrt{L_2 C_2})$.

\textbf{Coupling:} Mutual impedance $M = k \times \sqrt{L_1 L_2}$. The coupling coefficient $k$ (0--1) is determined by the transition geometry.

The transfer function $H(f) = I_2 / V_1$ shows two resonance peaks (at $f_1$ and $f_H$). Crucially: $k$ changes the \textbf{height} of the peaks (energy transfer = response), but not their \textbf{position} (frequency = filter action). Conversely, $C_2$ (chamber volume) shifts the $f_H$ peak without changing response.

\begin{tcolorbox}[colframe=keygreen, colback=green!5!white]
\textbf{Calculable:} With the physical values (mass, stiffness, volume, opening area) the transfer function is quantitatively determined. $k$ determines response, $C_2$ determines filter action --- these parameters are independent.
\end{tcolorbox}

\section{Valve Position --- Fine Adjustment}

In practice the valve opening can be offset by a few millimetres to at most one centimetre from the chamber end. The shift changes the effective acoustic neck length and thus $f_H$.

In the practical range (0--15\,mm), $f_H$ shifts by a few percent. Valve position is a fine adjustment --- the main lever is at the reed end (Port~1), where 70\,\% of conductance resides.

\section{Limits of the Model}

The two-port model shows relationships and directions. It does not say where the optimum lies.

\textbf{(a) Nonlinearity:} The reed switches periodically between small and large opening. The actual impedance modulation (Doc.~0013) is highly nonlinear.

\textbf{(b) Tonal target:} The ``optimal'' spectrum depends on the desired tonal character. The calculation shows the map; the ear walks the path.

\textbf{(c) Interactions:} Gap width (Doc.~0010), stiffness (Doc.~0012), chamber volume, valve opening and transition geometry all act together.

\begin{tcolorbox}[colframe=red!60!black, colback=red!5!white]
\textbf{This document is an explanatory model, not a design guide. Design decisions still require practical tests and experience.}
\end{tcolorbox}

\section{Where Does the Reed Couple? --- Centroid of Volume Flow}

The assumption that coupling occurs at the absolute end of the reed (tip) is a simplification. The reed is a cantilever --- its vibration shape follows the first Euler--Bernoulli bending mode.

Volume flow per unit length $\mathrm{d}Q/\mathrm{d}x \propto y(x)$ increases steeply towards the tip. The weighted \textbf{centroid of volume flow} lies at $x/L = 0{.}73$ --- not at the tip ($x/L = 1{.}0$). 50\,\% of total volume flow arises in the last 23\,\% of the reed length.

For a bass reed with $L = 70$\,mm the centroid lies at 51\,mm from the clamping point --- about 19\,mm before the tip.

\subsection{Impedance Mismatch}

The mechanical impedance of the reed and the acoustic impedance of the air differ by orders of magnitude. The reed is a \textbf{current generator} (high impedance) --- it drives a volume flow that the air accommodates.

$Z_\mathrm{reed}$ lies a factor 50--500 above $Z_\mathrm{air}$. The coupling is \textbf{not matched} in the sense of maximum power transfer ($Z_\mathrm{source} = Z_\mathrm{load}$). This is physically correct: the reed is self-excited and vibrates at its natural frequency, not at the chamber resonance. The air is ``carried along'', not ``driven''.

This is consistent with Doc.~0010: over 97\,\% of energy output occurs through air coupling, less than 3\,\% through the mechanical clamping.

\subsection{Consequence: Position of the Plate Edge}

When the reed plate covers the full reed length, the entire slot couples into the chamber. If the reed protrudes beyond the plate edge (deflection), the most energetic part --- the last 20--30\,\% of the reed --- bypasses the chamber and radiates directly into free space.

The position of the plate edge relative to the reed therefore determines what fraction of the volume flow is coupled into the chamber and what fraction radiates directly.

\section{Summary}

\begin{enumerate}
\item \textbf{Two-port system:} The chamber is open at both ends. Port~1 (reed) carries $\approx 70\,\%$ of conductance. The single-port model underestimates $f_H$.
\item \textbf{Fold invisible:} At 50\,Hz, the largest chamber dimension is 1.7\,\% of $\lambda$. Chamber shape is irrelevant --- only volume counts.
\item \textbf{Impedance matching:} The transition geometry acts as a transformer between bellows pressure (source) and reed (load). Maximum power transfer at $R_i = R_L$. Filter action ($f_H$) depends on the resonator, not the transformer.
\item \textbf{Equivalent circuit:} Reed and chamber form two coupled RLC resonators. $k$ determines response (energy transfer), $C_2 = V/(\rho c^2)$ determines filter action ($f_H$). These parameters are independent.
\item \textbf{Valve position = fine adjustment:} A few percent $f_H$ shift. The main lever is at the reed end.
\item \textbf{Centroid at $x/L = 0{.}73$:} Coupling does not occur at the reed tip, but $\approx 27\,\%$ before it. Plate edge position determines the coupling.
\item \textbf{$Z_\mathrm{reed} \gg Z_\mathrm{air}$:} The reed is a current generator. Impedances are not matched --- the air is carried along, not driven.
\item \textbf{Optimum is heard:} The calculation shows relationships and directions. Experience and ear training remain the key to success.
\end{enumerate}

\bigskip
\textit{The reed feeds, the chamber shapes, the valve radiates. The geometry at the transition determines the coupling.}


\end{document}
